\documentclass[a4paper,14pt]{extarticle}
\usepackage[utf8]{inputenc}
\usepackage[T2A]{fontenc}
\usepackage[english,russian]{babel}
\usepackage{geometry}
\usepackage{graphicx}
\usepackage{amsmath}
\usepackage{amsfonts}
\usepackage{amssymb}
\usepackage{hyperref}
\usepackage{caption}
\usepackage{subcaption}
\usepackage{listings}
\usepackage{xcolor}
\usepackage{float}

\geometry{left=30mm, right=15mm, top=20mm, bottom=20mm}

\definecolor{codegreen}{rgb}{0,0.6,0}
\definecolor{codegray}{rgb}{0.5,0.5,0.5}
\definecolor{codepurple}{rgb}{0.58,0,0.82}
\definecolor{backcolour}{rgb}{0.95,0.95,0.92}

\lstdefinestyle{mystyle}{
    backgroundcolor=\color{backcolour},   
    commentstyle=\color{codegreen},
    keywordstyle=\color{magenta},
    numberstyle=\tiny\color{codegray},
    stringstyle=\color{codepurple},
    basicstyle=\ttfamily\footnotesize,
    breakatwhitespace=false,         
    breaklines=true,                 
    captionpos=b,                    
    keepspaces=true,                 
    numbers=left,                    
    numbersep=5pt,                  
    showspaces=false,                
    showstringspaces=false,
    showtabs=false,                  
    tabsize=2,
    language=Python
}

\lstset{style=mystyle}

\begin{document}


\begin{titlepage}
\begin{center}
\large
Санкт-Петербургский политехнический университет \\
Петра Великого \\
Высшая школа прикладной математики и вычислительной физики \\
\vspace{3cm}

\textbf{Отчёт по лабораторной работе по дисциплине} \\
«Математическая статистика» \\
\vspace{0.5cm}

\textbf{ЛАБОРАТОРНАЯ РАБОТА} \\
«Гистограммы и плотности распределений» \\
\vspace{2cm}

Выполнил: Гуцалов Егор Никитич \\
Группа: 5030102/30201 \\
\vfill

Санкт-Петербург \\
2025
\end{center}
\end{titlepage}

\section*{2. Постановка задачи}

\textbf{Цель работы:} Знакомство с основными методами описательной статистики: визуализацией распределений, вычислением характеристик положения и рассеяния, анализом выбросов и построением эмпирических оценок функций распределения и плотности.

\textbf{Задачи:}
\begin{itemize}
    \item Сгенерировать выборки объемом $n = 10, 100, 1000$ для каждого из 5 распределений
    \item Построить на одном графике гистограмму выборки и теоретическую кривую плотности распределения
    \item Освоить принцип группировки данных и построения гистограмм
    \item Сделать вывод о влиянии размера выборки на определение характера распределения
\end{itemize}

\textbf{Исследуемые распределения:}
\begin{enumerate}
    \item Нормальное $N(0, 1)$
    \item Коши $C(0, 1)$
    \item Лапласа $L(0, \frac{1}{\sqrt{2}})$
    \item Пуассона $P(10)$
    \item Равномерное $U(-\sqrt{3}, \sqrt{3})$
\end{enumerate}

\section*{3. Теоретическая часть}

\subsection*{3.1 Гистограмма}

Гистограмма — это ступенчатая функция, которая аппроксимирует плотность распределения. Для построения гистограммы диапазон данных разбивается на интервалы (бины), и для каждого интервала подсчитывается количество попавших в него наблюдений.

Пусть имеется выборка $x_1, x_2, \ldots, x_n$. Диапазон разбивается на $k$ интервалов $[t_{j-1}, t_j)$, $j = 1, \ldots, k$. Тогда гистограмма плотности определяется как:

\begin{equation}
\hat{f}_n(x) = \frac{1}{n h_j} \sum_{i=1}^{n} I(x_i \in [t_{j-1}, t_j)), \quad x \in [t_{j-1}, t_j)
\end{equation}

где $h_j = t_j - t_{j-1}$ — ширина $j$-го интервала, $I(\cdot)$ — индикаторная функция.

\subsection*{3.2 Нормальное распределение}

Нормальное распределение $N(\mu, \sigma^2)$ задается плотностью:

\begin{equation}
f(x) = \frac{1}{\sigma\sqrt{2\pi}} e^{-\frac{(x-\mu)^2}{2\sigma^2}}
\end{equation}

где $\mu$ — математическое ожидание, $\sigma^2$ — дисперсия.

\subsection*{3.3 Распределение Коши}

Распределение Коши $C(x_0, \gamma)$ имеет плотность:

\begin{equation}
f(x) = \frac{1}{\pi \gamma \left[1 + \left(\frac{x-x_0}{\gamma}\right)^2\right]}
\end{equation}

где $x_0$ — параметр сдвига (медиана), $\gamma$ — параметр масштаба.

\subsection*{3.4 Распределение Лапласа}

Распределение Лапласа $L(\mu, b)$ задается плотностью:

\begin{equation}
f(x) = \frac{1}{2b} e^{-\frac{|x-\mu|}{b}}
\end{equation}

где $\mu$ — параметр сдвига, $b > 0$ — параметр масштаба.

\subsection*{3.5 Распределение Пуассона}

Распределение Пуассона $P(\lambda)$ — дискретное распределение:

\begin{equation}
P(X = k) = \frac{\lambda^k e^{-\lambda}}{k!}, \quad k = 0, 1, 2, \ldots
\end{equation}

где $\lambda > 0$ — интенсивность (среднее количество событий).

\subsection*{3.6 Равномерное распределение}

Равномерное распределение $U(a, b)$ имеет плотность:

\begin{equation}
f(x) = \begin{cases}
\frac{1}{b-a}, & a \leq x \leq b \\
0, & \text{иначе}
\end{cases}
\end{equation}

\section*{4. Реализация}

\textbf{Язык программирования:} Python 3.x

\textbf{Используемые библиотеки:}
\begin{itemize}
    \item \texttt{numpy} — генерация случайных выборок
    \item \texttt{matplotlib.pyplot} — визуализация данных
    \item \texttt{scipy.stats} — теоретические распределения
\end{itemize}

\textbf{Среда разработки:} Visual Studio Code

\subsection*{Основные этапы реализации:}

\begin{enumerate}
    \item Генерация выборок заданного объема из каждого распределения
    \item Построение гистограмм с автоматическим выбором количества бинов
    \item Наложение теоретических кривых плотности
    \item Сохранение графиков в файлы
\end{enumerate}

Пример кода генерации выборки и построения гистограммы:

\begin{lstlisting}[caption={Генерация выборки и построение гистограммы}]
import numpy as np
import matplotlib.pyplot as plt
from scipy import stats

n = 1000
sample = np.random.normal(loc=0, scale=1, size=n)

plt.hist(sample, bins='sqrt', density=True, 
         alpha=0.6, color='skyblue', 
         edgecolor='black', label='Histogram')

x = np.linspace(-4, 4, 1000)
pdf = stats.norm.pdf(x, loc=0, scale=1)
plt.plot(x, pdf, 'r-', linewidth=2, 
         label='Theoretical density')

plt.legend()
plt.grid(True, alpha=0.3)
plt.show()
\end{lstlisting}

\section*{5. Результаты}

\subsection*{5.1 Нормальное распределение $N(0, 1)$}

\begin{figure}[H]
\centering
\includegraphics[width=0.9\textwidth]{normal_distribution.png}
\caption{Гистограммы выборки и теоретическая плотность нормального распределения}
\label{fig:normal}
\end{figure}

\subsection*{5.2 Распределение Коши $C(0, 1)$}

\begin{figure}[H]
\centering
\includegraphics[width=0.9\textwidth]{cauchy_distribution.png}
\caption{Гистограммы выборки и теоретическая плотность распределения Коши}
\label{fig:cauchy}
\end{figure}

\subsection*{5.3 Распределение Лапласа $L(0, 1/\sqrt{2})$}

\begin{figure}[H]
\centering
\includegraphics[width=0.9\textwidth]{laplace_distribution.png}
\caption{Гистограммы выборки и теоретическая плотность распределения Лапласа}
\label{fig:laplace}
\end{figure}

\subsection*{5.4 Распределение Пуассона $P(10)$}

\begin{figure}[H]
\centering
\includegraphics[width=0.9\textwidth]{poisson_distribution.png}
\caption{Гистограммы выборки и теоретическая вероятность распределения Пуассона}
\label{fig:poisson}
\end{figure}

\subsection*{5.5 Равномерное распределение $U(-\sqrt{3}, \sqrt{3})$}

\begin{figure}[H]
\centering
\includegraphics[width=0.9\textwidth]{uniform_distribution.png}
\caption{Гистограммы выборки и теоретическая плотность равномерного распределения}
\label{fig:uniform}
\end{figure}

\section*{6. Обсуждение}

В ходе выполнения лабораторной работы были получены следующие результаты:

\begin{enumerate}
    \item \textbf{Влияние объема выборки:} При увеличении объема выборки от $n=10$ до $n=1000$ гистограмма все лучше приближается к теоретической кривой плотности. Это согласуется с законом больших чисел.
    
    \item \textbf{Нормальное распределение:} Гистограмма при $n=1000$ практически идеально совпадает с теоретической кривой. Распределение симметрично относительно нуля.
    
    \item \textbf{Распределение Коши:} Имеет более тяжелые хвосты по сравнению с нормальным распределением. Наблюдаются значительные выбросы, что характерно для данного распределения.
    
    \item \textbf{Распределение Лапласа:} Отличается более острым пиком в центре по сравнению с нормальным распределением.
    
    \item \textbf{Распределение Пуассона:} Как дискретное распределение, визуализировано с помощью stem-plot. При $\lambda=10$ приближается к нормальному.
    
    \item \textbf{Равномерное распределение:} Гистограмма имеет характерную прямоугольную форму, все значения в интервале равновероятны.
\end{enumerate}

\section*{7. Выводы}

\begin{enumerate}
    \item Освоен принцип построения гистограмм для визуализации эмпирических распределений.
    
    \item Показано, что с ростом объема выборки эмпирическое распределение стремится к теоретическому.
    
    \item Изучены характеристики пяти различных типов распределений: непрерывных (нормальное, Коши, Лапласа, равномерное) и дискретного (Пуассона).
    
    \item Подтверждено, что выбор оптимального количества интервалов группировки важен для адекватного представления данных.
    
    \item Приобретены практические навыки работы с библиотеками Python для статистического анализа данных.
\end{enumerate}

\section*{8. Список литературы}

\begin{enumerate}
    \item Гмурман В.Е. Теория вероятностей и математическая статистика. — М.: Высшая школа, 2003.
    
    \item Кобецкая И.А. Математическая статистика: Учебное пособие. — СПб.: Лань, 2019.
    
    \item McKinney W. Python for Data Analysis. — O'Reilly Media, 2022.
    
    \item SciPy documentation. URL: \url{https://docs.scipy.org/doc/scipy/}
    
    \item Matplotlib documentation. URL: \url{https://matplotlib.org/stable/contents.html}
\end{enumerate}

\section*{9. Приложение}

Исходный код программы размещен в репозитории GitHub:

\url{https://github.com/Ftp-Glich/Statistics}

\end{document}